% ============================================================
% ab-vorlage.tex – gemeinsame Vorlage der Arbeitsblätter
% englisch/13/questions-on-the-text/arbeitsblaetter/
% Übersetzen mit:  xelatex <datei>.tex   (zweimal, wegen Seitenzahlen)
% Lösungsfassung:  Einzeiler-Datei <datei>-lsg.tex mit
%                  \def\mitloesung{}\input{<datei>}
% Schriften: Carlito (Fließtext), Poppins (Überschriften), TeX Gyre Pagella (Texte)
% ============================================================
\documentclass[10pt,a4paper]{article}
\usepackage[a4paper,top=12mm,bottom=13mm,left=14mm,right=14mm,headsep=4mm,headheight=12pt,footskip=7mm]{geometry}
\usepackage{fontspec}
\setmainfont{Carlito}
\newfontfamily\headfont{Poppins}[UprightFont=* Medium, BoldFont=* Medium]
\newfontfamily\headlight{Poppins Light}
\newfontfamily\lorafont{TeX Gyre Pagella}
\usepackage{xcolor}
\definecolor{prim}{HTML}{2563EB}
\definecolor{primdark}{HTML}{1E3A8A}
\definecolor{sec}{HTML}{7C3AED}
\definecolor{primtint}{HTML}{EFF6FF}
\definecolor{sectint}{HTML}{F5F3FF}
\definecolor{ink}{HTML}{1E293B}
\definecolor{muted}{HTML}{64748B}
\definecolor{rule}{HTML}{CBD5E1}
\definecolor{good}{HTML}{059669}
\definecolor{goodtint}{HTML}{ECFDF5}
\definecolor{warn}{HTML}{B45309}
\definecolor{warntint}{HTML}{FFFBEB}
\definecolor{bad}{HTML}{B91C1C}
\definecolor{badtint}{HTML}{FEF2F2}
\color{ink}
\usepackage[most]{tcolorbox}
\usepackage{tikz}
\usetikzlibrary{positioning,calc,arrows.meta,shapes.misc}
\usepackage{enumitem}
\setlist{nosep,leftmargin=1.1em,topsep=1pt}
\setlist[itemize]{label=\textcolor{prim}{\small\textbullet}}
\usepackage{tabularx,array,booktabs,colortbl}
\usepackage{multicol}
\usepackage{fancyhdr}
\usepackage{lastpage}
\usepackage{amssymb}
\usepackage{microtype}
\setlength{\parindent}{0pt}
\setlength{\parskip}{0pt}
\linespread{1.02}

% ---------- Kopf- und Fußzeile ----------
\newcommand{\wskopf}{English · Q13 · Science and Technology}
\pagestyle{fancy}
\fancyhf{}
\renewcommand{\headrulewidth}{0pt}
\fancyhead[L]{\footnotesize\headlight\textcolor{muted}{\wskopf}}
\fancyhead[R]{\footnotesize\headlight\textcolor{muted}{Name: \rule{38mm}{0.3pt}\quad Date: \rule{18mm}{0.3pt}}}
\fancyfoot[L]{\scriptsize\textcolor{muted}{edu-mrh.de/englisch/13/questions-on-the-text}}
\fancyfoot[R]{\scriptsize\textcolor{muted}{\thepage\,/\,\pageref*{LastPage}}}

% ---------- Titel ----------
% \wstitel{Kürzel}{Titel}{Untertitel}
\newcommand{\wstitel}[3]{%
  \begin{tcolorbox}[enhanced,colback=prim,colframe=prim,arc=2mm,boxrule=0pt,
     left=4mm,right=4mm,top=2.2mm,bottom=2.2mm,
     interior style={left color=primdark,right color=sec}]
    \color{white}{\headlight\footnotesize #1}\\[0.3mm]
    {\headfont\LARGE #2}\\[0.6mm]
    {\small #3}
  \end{tcolorbox}\vspace{1mm}}

% ---------- Boxen ----------
% \begin{wsbox}{Nummer}{Überschrift} ... \end{wsbox}
\newtcolorbox{wsbox}[3][]{enhanced,breakable,colback=white,colframe=rule,boxrule=0.5pt,arc=1.5mm,
  left=2.5mm,right=2.5mm,top=1.3mm,bottom=1.3mm,
  title={\headfont\normalsize\textcolor{white}{\colorbox{prim}{\,#2\,}}\;\textcolor{ink}{#3}},
  coltitle=ink,colbacktitle=white,attach title to upper={\par\vspace{0.6mm}},before skip=1.4mm,after skip=1.4mm,#1}
\newtcolorbox{tintbox}[1][]{enhanced,colback=primtint,colframe=primtint,boxrule=0pt,arc=1.2mm,
  left=2mm,right=2mm,top=1.2mm,bottom=1.2mm,#1}
\newtcolorbox{secbox}[1][]{enhanced,colback=sectint,colframe=sectint,boxrule=0pt,arc=1.2mm,
  left=2mm,right=2mm,top=1.2mm,bottom=1.2mm,#1}
\newtcolorbox{goodbox}[1][]{enhanced,colback=goodtint,colframe=goodtint,boxrule=0pt,arc=1.2mm,
  left=2mm,right=2mm,top=1.2mm,bottom=1.2mm,#1}
\newtcolorbox{warnbox}[1][]{enhanced,colback=warntint,colframe=warntint,boxrule=0pt,arc=1.2mm,
  left=2mm,right=2mm,top=1.2mm,bottom=1.2mm,#1}
\newtcolorbox{badbox}[1][]{enhanced,colback=badtint,colframe=badtint,boxrule=0pt,arc=1.2mm,
  left=2mm,right=2mm,top=1.2mm,bottom=1.2mm,#1}

\newcommand{\kl}[1]{{\headfont\small\textcolor{prim}{#1}}}      % kleine Zwischenüberschrift
\newcommand{\klsec}[1]{{\headfont\small\textcolor{sec}{#1}}}
\newcommand{\plus}{\textcolor{good}{\textbf{+}}\ }
\newcommand{\minus}{\textcolor{bad}{\textbf{–}}\ }
\newcommand{\pfeil}{\textcolor{prim}{$\rightarrow$}\ }
\newcommand{\linie}[1][\linewidth]{\rule{#1}{0.3pt}}
% ---------- Lösungsschalter ----------
% \lsg[Breite]{Text}        Lücke mit Unterstrich; Lösung erscheint nur mit \mitloesung
% \lsglinien{n}{Text}       n Schreiblinien; Lösung wird auf die Linien geschrieben
% Beide reservieren in beiden Fassungen exakt denselben Platz.
\newif\ifloesung
\ifdefined\mitloesung \loesungtrue \else \loesungfalse \fi
\newcommand{\lsgstil}{\color{sec}\itshape}
\newcommand{\lsg}[2][30mm]{%
  \makebox[#1][l]{\rlap{\raisebox{-0.8mm}{\textcolor{rule}{\rule{#1}{0.4pt}}}}%
  \ifloesung{\lsgstil\small #2}\fi}}
\newlength{\lsgzeile}\setlength{\lsgzeile}{7mm}
\newcommand{\lsglinien}[2]{%
  \par\noindent
  \begin{tikzpicture}[x=1mm,y=1mm]
    \pgfmathsetmacro{\hoehe}{#1*7}
    \path[use as bounding box] (0,0.5) rectangle (\linewidth,-\hoehe);
    \foreach \i in {1,...,#1} {\draw[rule,line width=0.4pt] (0,-\i*7) -- (\linewidth,-\i*7);}
    \ifloesung
      \node[anchor=north west,inner sep=0pt,text width=\linewidth,align=left,
            font=\lsgstil\small] at (0,-3.3) {\setlength{\baselineskip}{\lsgzeile}#2\par};
    \fi
  \end{tikzpicture}\par}
\newcommand{\schreiblinien}[1]{\lsglinien{#1}{}}
